\documentclass[10pt,a4paper]{amsart}
\usepackage{amsmath,amssymb,amsthm,graphicx}
% Retained preamble for the Stacks Project chapter snapshot in this directory.
% The chapter's own preamble.tex is a copy of the shared Stacks Project preamble; the block
% below is sliced MECHANICALLY, line for line, from it. Nothing was added, removed or
% redefined; the wrapper class, the xr-hyper/hyperref wiring and the externaldocument list
% are omitted because the assembler supplies the document class.
%
% lines taken from preamble.tex: 3, 5, 9, 11, 12, 156, 157, 158, ... (51 lines total)
%
% package coverage restored from the chapter's own preamble (2026-10-13): xy, xy-2cell

\usepackage[all]{xy}
\xyoption{2cell}
\UseAllTwocells

\newenvironment{reference}{\comment}{\endcomment}
\newenvironment{slogan}{\comment}{\endcomment}
\newenvironment{history}{\comment}{\endcomment}
\theoremstyle{plain}
\newtheorem{theorem}[subsection]{Theorem}
\newtheorem{proposition}[subsection]{Proposition}
\newtheorem{lemma}[subsection]{Lemma}
\theoremstyle{definition}
\newtheorem{definition}[subsection]{Definition}
\newtheorem{example}[subsection]{Example}
\newtheorem{exercise}[subsection]{Exercise}
\newtheorem{situation}[subsection]{Situation}
\theoremstyle{remark}
\newtheorem{remark}[subsection]{Remark}
\newtheorem{remarks}[subsection]{Remarks}
\numberwithin{equation}{subsection}
\def\lim{\mathop{\mathrm{lim}}\nolimits}
\def\colim{\mathop{\mathrm{colim}}\nolimits}
\def\Spec{\mathop{\mathrm{Spec}}}
\def\Hom{\mathop{\mathrm{Hom}}\nolimits}
\def\Ext{\mathop{\mathrm{Ext}}\nolimits}
\def\SheafHom{\mathop{\mathcal{H}\!\mathit{om}}\nolimits}
\def\SheafExt{\mathop{\mathcal{E}\!\mathit{xt}}\nolimits}
\def\Sch{\mathit{Sch}}
\def\Mor{\mathop{\mathrm{Mor}}\nolimits}
\def\Ob{\mathop{\mathrm{Ob}}\nolimits}
\def\Sh{\mathop{\mathit{Sh}}\nolimits}
\def\NL{\mathop{N\!L}\nolimits}
\def\CH{\mathop{\mathrm{CH}}\nolimits}
\def\proetale{{pro\text{-}\acute{e}tale}}
\def\etale{{\acute{e}tale}}
\def\QCoh{\mathit{QCoh}}
\def\Ker{\mathop{\mathrm{Ker}}}
\def\Im{\mathop{\mathrm{Im}}}
\def\Coker{\mathop{\mathrm{Coker}}}
\def\Coim{\mathop{\mathrm{Coim}}}
\DeclareMathSymbol{\boxtimes}{\mathbin}{AMSa}{"02}
\def\QCohstack{\mathcal{QC}\!\mathit{oh}}
\def\Cohstack{\mathcal{C}\!\mathit{oh}}
\def\Spacesstack{\mathcal{S}\!\mathit{paces}}
\def\Quotfunctor{\mathrm{Quot}}
\def\Hilbfunctor{\mathrm{Hilb}}
\def\Curvesstack{\mathcal{C}\!\mathit{urves}}
\def\Polarizedstack{\mathcal{P}\!\mathit{olarized}}
\def\Complexesstack{\mathcal{C}\!\mathit{omplexes}}
\def\Pic{\mathop{\mathrm{Pic}}\nolimits}
\def\Picardstack{\mathcal{P}\!\mathit{ic}}
\def\Picardfunctor{\mathrm{Pic}}
\def\Deformationcategory{\mathcal{D}\!\mathit{ef}}
\usepackage{cleveref}
\crefname{originallemma}{original lemma}{original lemmas}
\Crefname{originallemma}{Original lemma}{Original lemmas}
\makeatletter
\newlabel{lemma-Mor-s-lfp}{{0DPN}{source}{}{sourcepointer-1}{}}
\newlabel{lemma-Mor-s-lfp@cref}{{[originallemma][0][]0DPN}{[1][0][]source}{}{}{}}
\makeatother
\crefname{originallemma}{original lemma}{original lemmas}
\Crefname{originallemma}{Original lemma}{Original lemmas}
\makeatletter
\newlabel{lemma-hilb-qc-over-base}{{0DPI}{source}{}{sourcepointer-2}{}}
\newlabel{lemma-hilb-qc-over-base@cref}{{[originallemma][0][]0DPI}{[1][0][]source}{}{}{}}
\makeatother
\crefname{originallemma}{original lemma}{original lemmas}
\Crefname{originallemma}{Original lemma}{Original lemmas}
\makeatletter
\newlabel{limits-lemma-descend-finite-presentation}{{01ZM}{source}{}{sourcepointer-3}{}}
\newlabel{limits-lemma-descend-finite-presentation@cref}{{[originallemma][0][]01ZM}{[1][0][]source}{}{}{}}
\makeatother
\crefname{originallemma}{original lemma}{original lemmas}
\Crefname{originallemma}{Original lemma}{Original lemmas}
\makeatletter
\newlabel{limits-lemma-descend-invertible-modules}{{0B8W}{source}{}{sourcepointer-4}{}}
\newlabel{limits-lemma-descend-invertible-modules@cref}{{[originallemma][0][]0B8W}{[1][0][]source}{}{}{}}
\makeatother
\crefname{originallemma}{original lemma}{original lemmas}
\Crefname{originallemma}{Original lemma}{Original lemmas}
\makeatletter
\newlabel{limits-lemma-descend-separated-finite-presentation}{{01ZQ}{source}{}{sourcepointer-5}{}}
\newlabel{limits-lemma-descend-separated-finite-presentation@cref}{{[originallemma][0][]01ZQ}{[1][0][]source}{}{}{}}
\makeatother
\crefname{originallemma}{original lemma}{original lemmas}
\Crefname{originallemma}{Original lemma}{Original lemmas}
\makeatletter
\newlabel{limits-lemma-eventually-proper}{{081F}{source}{}{sourcepointer-6}{}}
\newlabel{limits-lemma-eventually-proper@cref}{{[originallemma][0][]081F}{[1][0][]source}{}{}{}}
\makeatother
\crefname{originallemma}{original lemma}{original lemmas}
\Crefname{originallemma}{Original lemma}{Original lemmas}
\makeatletter
\newlabel{limits-lemma-limit-ample}{{09MT}{source}{}{sourcepointer-7}{}}
\newlabel{limits-lemma-limit-ample@cref}{{[originallemma][0][]09MT}{[1][0][]source}{}{}{}}
\makeatother
\crefname{originallemma}{original lemma}{original lemmas}
\Crefname{originallemma}{Original lemma}{Original lemmas}
\makeatletter
\newlabel{properties-lemma-ample-on-product}{{0DNK}{source}{}{sourcepointer-8}{}}
\newlabel{properties-lemma-ample-on-product@cref}{{[originallemma][0][]0DNK}{[1][0][]source}{}{}{}}
\makeatother
\crefname{originallemma}{original lemma}{original lemmas}
\Crefname{originallemma}{Original lemma}{Original lemmas}
\makeatletter
\newlabel{remark-Mor-numerical}{{0DPQ}{source}{}{sourcepointer-9}{}}
\newlabel{remark-Mor-numerical@cref}{{[originallemma][0][]0DPQ}{[1][0][]source}{}{}{}}
\makeatother
\crefname{originallemma}{original lemma}{original lemmas}
\Crefname{originallemma}{Original lemma}{Original lemmas}
\makeatletter
\newlabel{spaces-divisors-definition-relatively-ample}{{0D31}{source}{}{sourcepointer-10}{}}
\newlabel{spaces-divisors-definition-relatively-ample@cref}{{[originallemma][0][]0D31}{[1][0][]source}{}{}{}}
\makeatother
\crefname{originallemma}{original lemma}{original lemmas}
\Crefname{originallemma}{Original lemma}{Original lemmas}
\makeatletter
\newlabel{spaces-morphisms-lemma-quasi-compact-local}{{03KG}{source}{}{sourcepointer-11}{}}
\newlabel{spaces-morphisms-lemma-quasi-compact-local@cref}{{[originallemma][0][]03KG}{[1][0][]source}{}{}{}}
\makeatother
\title{The polynomial piece of the Mor space is separated of finite presentation}
\author{The Stacks Project Authors}
\date{stacks-project:moduli.tex:local-snapshot-20261012; GNU FDL 1.2 or later}
\begin{document}
\maketitle
\noindent\textit{Editorial note.} Statement span is lines 1898-1911 of sources/stacks-moduli/pinned-original-source.tex; proof span is lines 1913-1962, the complete author proof (48 body lines).\par\medskip
\noindent\textit{Editorial note.} Navigation stubs: 11 cited-but-not-included labels are supplied as non-load-bearing pointers with their pinned tags where the snapshot tag index has them.\par\medskip
\noindent\textit{Editorial note.} extra\_packages is empty on purpose: the retained preamble already loads xy with options and a second declaration clashes in the assembler wrapper.\par\medskip
\noindent\textit{Editorial note.} Text is retained verbatim, including the author's own typos and placeholders.\par\medskip
\noindent\textit{Editorial note.} Queue index 7 of queue\_moduli.json (pb=2158, stmt 1898-1911, proof 1913-1962).\par\medskip
\noindent\textit{Editorial note.} stacks-moduli wave-1: fresh adjudication accept (2026-10-14); pool per Chief order 08:5x (dpa superseded to lane-s9); moduli.tex 53 labels/48 with proof judged 25 accept / 23 reject / 5 pb=0 env-dup rows excluded; spans cite-free and reference/slogan/history-free so no patchenv needed.\par\medskip
\noindent\textit{Allowed standard prerequisites (editorial).}\begin{itemize}
\item Explicit prerequisite boundary. The cited inputs are prerequisites; the argument in the retained span is not replaced by any cited result.
\end{itemize}\medskip
\ifdefined\hypertarget\hypertarget{sourcepointer-1}{}\fi
\noindent\textit{Source reference.} Original lemma 0DPN: sources/stacks-moduli/pinned-original-source.tex, source lines 1843--1846. Cited by the retained statement or proof but not part of the collected spans; supplied only as a navigation pointer. The argument does not depend on its content.\par\medskip
\ifdefined\hypertarget\hypertarget{sourcepointer-2}{}\fi
\noindent\textit{Source reference.} Original lemma 0DPI: sources/stacks-moduli/pinned-original-source.tex, source lines 1476--1486. Cited by the retained statement or proof but not part of the collected spans; supplied only as a navigation pointer. The argument does not depend on its content.\par\medskip
\ifdefined\hypertarget\hypertarget{sourcepointer-3}{}\fi
\noindent\textit{Source reference.} Original lemma 01ZM: sources/stacks-moduli/pinned-original-source.tex, source lines 1962--1962. Cited by the retained statement or proof but not part of the collected spans; supplied only as a navigation pointer. The argument does not depend on its content.\par\medskip
\ifdefined\hypertarget\hypertarget{sourcepointer-4}{}\fi
\noindent\textit{Source reference.} Original lemma 0B8W: sources/stacks-moduli/pinned-original-source.tex, source lines 1962--1962. Cited by the retained statement or proof but not part of the collected spans; supplied only as a navigation pointer. The argument does not depend on its content.\par\medskip
\ifdefined\hypertarget\hypertarget{sourcepointer-5}{}\fi
\noindent\textit{Source reference.} Original lemma 01ZQ: sources/stacks-moduli/pinned-original-source.tex, source lines 1962--1962. Cited by the retained statement or proof but not part of the collected spans; supplied only as a navigation pointer. The argument does not depend on its content.\par\medskip
\ifdefined\hypertarget\hypertarget{sourcepointer-6}{}\fi
\noindent\textit{Source reference.} Original lemma 081F: sources/stacks-moduli/pinned-original-source.tex, source lines 1962--1962. Cited by the retained statement or proof but not part of the collected spans; supplied only as a navigation pointer. The argument does not depend on its content.\par\medskip
\ifdefined\hypertarget\hypertarget{sourcepointer-7}{}\fi
\noindent\textit{Source reference.} Original lemma 09MT: sources/stacks-moduli/pinned-original-source.tex, source lines 1962--1962. Cited by the retained statement or proof but not part of the collected spans; supplied only as a navigation pointer. The argument does not depend on its content.\par\medskip
\ifdefined\hypertarget\hypertarget{sourcepointer-8}{}\fi
\noindent\textit{Source reference.} Original lemma 0DNK: sources/stacks-moduli/pinned-original-source.tex, source lines 1962--1962. Cited by the retained statement or proof but not part of the collected spans; supplied only as a navigation pointer. The argument does not depend on its content.\par\medskip
\ifdefined\hypertarget\hypertarget{sourcepointer-9}{}\fi
\noindent\textit{Source reference.} Original lemma 0DPQ: sources/stacks-moduli/pinned-original-source.tex, source lines 1870--1896. Cited by the retained statement or proof but not part of the collected spans; supplied only as a navigation pointer. The argument does not depend on its content.\par\medskip
\ifdefined\hypertarget\hypertarget{sourcepointer-10}{}\fi
\noindent\textit{Source reference.} Original lemma 0D31: sources/stacks-moduli/pinned-original-source.tex, source lines 1962--1962. Cited by the retained statement or proof but not part of the collected spans; supplied only as a navigation pointer. The argument does not depend on its content.\par\medskip
\ifdefined\hypertarget\hypertarget{sourcepointer-11}{}\fi
\noindent\textit{Source reference.} Original lemma 03KG: sources/stacks-moduli/pinned-original-source.tex, source lines 1962--1962. Cited by the retained statement or proof but not part of the collected spans; supplied only as a navigation pointer. The argument does not depend on its content.\par\medskip
\begin{lemma}
\label{lemma-Mor-qc-over-base}
With $B, X, Y$ as in the introduction of this section, let
$\mathcal{L}$ be ample on $X/B$ and let $\mathcal{N}$ be ample on $Y/B$.
See Divisors on Spaces, Definition
\ref{spaces-divisors-definition-relatively-ample}.
Let $P$ be a numerical polynomial. Then
$$
\mathit{Mor}^{P, \mathcal{M}}_B(Y, X) \longrightarrow B
$$
is separated and of finite presentation where
$\mathcal{M} = \text{pr}_1^*\mathcal{N}
\otimes_{\mathcal{O}_{Y \times_B X}} \text{pr}_2^*\mathcal{L}$.
\end{lemma}

\begin{proof}
By Lemma \ref{lemma-Mor-s-lfp} the morphism $\mathit{Mor}_B(Y, X) \to B$
is separated and locally of finite presentation. Thus it suffices to
show that the open and closed subspace $\mathit{Mor}^{P, \mathcal{M}}_B(Y, X)$
of Remark \ref{remark-Mor-numerical} is quasi-compact over $B$.

\medskip\noindent
The question is \'etale local on $B$
(Morphisms of Spaces, Lemma \ref{spaces-morphisms-lemma-quasi-compact-local}).
Thus we may assume $B$ is affine.

\medskip\noindent
Assume $B = \Spec(\Lambda)$. Note that $X$ and
$Y$ are schemes and that $\mathcal{L}$ and $\mathcal{N}$ are ample
invertible sheaves on $X$ and $Y$ (this follows immediately from the
definitions). Write $\Lambda = \colim \Lambda_i$ as the
colimit of its finite type $\mathbf{Z}$-subalgebras. Then
we can find an $i$ and a system $X_i, Y_i, \mathcal{L}_i, \mathcal{N}_i$
as in the lemma over $B_i = \Spec(\Lambda_i)$ whose base change to
$B$ gives $X, Y, \mathcal{L}, \mathcal{N}$. This follows from
Limits, Lemmas
\ref{limits-lemma-descend-finite-presentation} (to find $X_i$, $Y_i$),
\ref{limits-lemma-descend-invertible-modules} (to find $\mathcal{L}_i$,
$\mathcal{N}_i$), \ref{limits-lemma-descend-separated-finite-presentation}
(to make $X_i \to B_i$ separated), \ref{limits-lemma-eventually-proper}
(to make $Y_i \to B_i$ proper), and \ref{limits-lemma-limit-ample}
(to make $\mathcal{L}_i$, $\mathcal{N}_i$ ample).
Because
$$
\mathit{Mor}_B(Y, X) = B \times_{B_i} \mathit{Mor}_{B_i}(Y_i, X_i)
$$
and similarly for $\mathit{Mor}^P_B(Y, X)$ we reduce
to the case discussed in the next paragraph.

\medskip\noindent
Assume $B$ is a Noetherian affine scheme. By
Properties, Lemma \ref{properties-lemma-ample-on-product}
we see that $\mathcal{M}$ is ample. By Lemma \ref{lemma-hilb-qc-over-base}
we see that $\Hilbfunctor^{P, \mathcal{M}}_{Y \times_B X/B}$ is of
finite presentation over $B$ and hence Noetherian.
By construction
$$
\mathit{Mor}^{P, \mathcal{M}}_B(Y, X) =
\mathit{Mor}_B(Y, X) \cap
\Hilbfunctor^{P, \mathcal{M}}_{Y \times_B X/B}
$$
is an open subspace of $\Hilbfunctor^{P, \mathcal{M}}_{Y \times_B X/B}$ and
hence quasi-compact (as an open of a Noetherian algebraic space
is quasi-compact).
\end{proof}

\end{document}
